\subsection{Interpretation of the Main Result}

The final test results produce different rankings across the
three loss functions. TVP-HAR-X achieves the lowest MAE and
RMSE on the log realized-variance scale, whereas Fixed HAR-X
achieves the lowest QLIKE on the original variance scale.

MAE and RMSE evaluate errors in
\(y_t=\log(RV_t)\). MAE summarizes the typical absolute
deviation, while RMSE assigns greater weight to large
log-scale errors.

QLIKE evaluates the implied variance forecast
\(\widehat{RV}_t=\exp(\widehat{y}_t)\). Its loss depends on the
ratio \(RV_t/\widehat{RV}_t\), making it particularly sensitive
to substantial underprediction during high-volatility periods
\citep{patton2011}. A model may therefore reduce average
log-scale errors while recording a higher QLIKE because of a
small number of unfavorable variance forecasts.

The test results suggest that coefficient adaptation improves
the typical accuracy of log-volatility forecasts. Fixed HAR-X,
however, provides slightly better performance under the
variance-scale loss. The empirical value of parameter
variation consequently depends on the forecasting objective.

\subsection{Validation and Test Performance}

TVP-HAR-X reduces validation QLIKE by approximately \(4.30\%\)
relative to Fixed HAR-X, but its test QLIKE is \(0.76\%\)
higher. This change in ranking illustrates the importance of
retaining a separate final test period.

Fixed HAR-X estimates a single coefficient vector from the
development sample and produces relatively stable forecasts.
TVP-HAR-X introduces gradual coefficient movement, reducing
bias when predictive relationships change while adding
state-estimation uncertainty.

The relative performance of the two specifications can be
understood as a trade-off between adaptation and estimation
uncertainty. During validation, parameter variation improves
all three losses. During the test period, it continues to
reduce MAE and RMSE, but the gain is insufficient to lower
QLIKE.

The difference across years also indicates that forecast-model
rankings are sensitive to market conditions. Chronological
validation is therefore more informative than a random split
for this application.

\subsection{Incremental Value of the HAR-X Information Set}

The most substantial model improvement occurs between AR(24)
and Fixed HAR-X. On the final test sample, Fixed HAR-X reduces
QLIKE by approximately \(5.75\%\) relative to AR(24).

AR(24) is a demanding benchmark because it estimates a
separate coefficient for each of the previous 24 hourly
volatility observations. The HAR-X improvement therefore
provides evidence that the structured predictor set contains
information beyond an unrestricted one-day lag specification.

The HAR variables extend the information set to the preceding
week without requiring 168 individual lag coefficients. This
multi-horizon representation follows the HAR framework of
\citet{corsi2009} and provides a parsimonious description of
persistent volatility.

HAR-X also incorporates current market activity. Among the
additional variables, log trade count displays the strongest
and most stable association with next-hour volatility. The
forecasting information set thus combines recent volatility,
longer-horizon volatility conditions, and current trading
activity.

\subsection{Interpretation of Trading Activity}

Trading activity is the most informative additional predictor
in both HAR-X specifications. In the fixed model, its
standardized coefficient is \(0.3379\), with a HAC
\(p\)-value below \(0.001\). During validation, its
time-varying coefficient remains positive and ranges from
approximately \(0.441\) to \(0.651\).

These estimates indicate that unusually active trading hours
tend to be followed by higher next-hour log realized variance,
conditional on volatility measured at the four HAR horizons.

The result is consistent with the interpretation of trade
count as a proxy for information arrival and market attention.
It is related to the mixture-of-distributions framework of
\citet{clark1973} and the transaction-count evidence reported
by \citet{jones1994}. Increased activity may reflect periods
in which market participants are processing new information
and adjusting positions more intensively.

The estimated coefficient represents a conditional predictive
association. Trading activity and volatility may respond
jointly to latent information or broader changes in market
conditions, so the model does not identify a causal effect.

\subsection{Downside Semivariance}

Downside semivariance is statistically insignificant in the
fixed HAR-X model, and its time-varying coefficient changes
sign during validation.

Realized semivariance separates the contributions of negative
and positive intraday returns to total variation
\citep{barndorffnielsen2010}. In the present specification,
total realized variance and the current-hour HAR component
already contain the squared magnitudes of both positive and
negative returns. Downside semivariance must therefore provide
information beyond the overall volatility level to obtain a
distinct effect.

The estimates suggest that this additional contribution is
small on average at the one-hour forecast horizon. The
logarithmic transformation may also reduce the influence of a
limited number of extreme downside observations. The
time-varying estimates do not reveal a persistent positive or
negative relationship.

\subsection{Aggressive-Buying Proxy}

The aggressive-buying proxy has a coefficient close to zero in
the fixed model and changes sign in the TVP specification.

The signed measure records the balance between taker-initiated
buying and selling. Large imbalances in either direction may
coincide with elevated volatility, which can weaken the average
linear effect of a signed variable.

An absolute or nonlinear imbalance measure could capture this
possibility more directly. Such alternatives were not included
in the pre-specified model and would require evaluation on a
new holdout sample.

The weak result may also reflect the information available in
candlestick data. Aggregated taker-buy volume does not capture
order-book depth, bid--ask spreads, cancellations, or
trade-level price impact.

\subsection{Time-Varying Coefficient Patterns}

The state estimates reveal substantial differences in
coefficient stability across predictors.

The six-hour and twenty-four-hour volatility coefficients
remain positive and vary within relatively narrow ranges,
indicating stable short- and medium-horizon persistence during
validation.

The weekly coefficient fluctuates around zero, suggesting that
the incremental contribution of longer-horizon volatility
changes across market conditions. The downside-semivariance
and aggressive-buying coefficients also cross zero and display
limited average effects.

Trading activity remains positive throughout validation,
although its magnitude changes over time. Its predictive
contribution is therefore persistent within this period but
not constant.

The state paths provide evidence that some components of the
forecasting relationship are more stable than others.
Short- and medium-horizon volatility persistence is relatively
stable, while weekly and directional effects display greater
variation. These patterns supplement the forecast comparison
by showing how the conditional relationships evolve over time.

\subsection{Stationarity and Parameter Variation}

The ADF and KPSS results provide complementary evidence on the
behavior of hourly log realized variance. The ADF test rejects
a unit root, supporting the mean-reverting structure captured
by AR(24). The KPSS test rejects global level stationarity
\citep{kpss1992}, indicating that the level or dependence
structure may vary over longer periods.

This evidence is consistent with previous findings that
structural breaks affect Bitcoin volatility forecasting
\citep{shen2020}. It also provides a statistical motivation
for allowing gradual coefficient variation in the HAR-X
relationship.

TVP-HAR-X retains the mean-reverting dependence represented by
the HAR components while permitting their coefficients to
evolve. The state-space specification therefore accommodates
parameter instability within an otherwise persistent
volatility process.

\subsection{Residual Dependence}

Ljung--Box portmanteau tests reject residual white noise for
AR(24), Fixed HAR-X, and TVP-HAR-X
\citep{ljungbox1978}. None of the specifications fully captures
the temporal dependence in hourly volatility.

The remaining correlation may reflect long-memory behavior,
intraday seasonality, nonlinear responses, omitted market
variables, or changes in conditional variance. However,
first-order residual and innovation autocorrelations are close
to zero. For TVP-HAR-X, the lag-one autocorrelation of the
standardized innovations is approximately \(0.0007\).

Given the large diagnostic samples, small correlations
distributed across several lags can produce statistically
significant portmanteau statistics. The residual diagnostics
therefore indicate model incompleteness, while the
out-of-sample results remain informative for relative forecast
comparison.

\subsection{Statistical and Practical Significance}

The HAC loss-differential tests show that TVP-HAR-X reduces
absolute and squared error relative to Fixed HAR-X. The
improvements are \(1.33\%\) for absolute error and \(1.87\%\)
for squared error, with one-sided \(p\)-values below \(0.001\).

These gains are statistically precise but modest in magnitude.
Their practical relevance depends on the application. Small
improvements may accumulate when forecasts are generated
hourly and used repeatedly for risk limits, margin
calculations, or position sizing. In settings where
implementation simplicity is more important, the additional
estimation and monitoring required by the state-space model may
outweigh the reduction in average error.

Fixed HAR-X is simpler to estimate and deploy, while
TVP-HAR-X provides lower log-scale forecast errors and
information about changes in predictor importance. Model choice
therefore depends on the relevant loss function and the value
assigned to adaptability.

\subsection{Empirical Design}

Several features of the design strengthen the forecast
comparison. The sample is divided chronologically, and all
preprocessing parameters are estimated from the training
period. Model selection is completed using validation data
before the final test sample is evaluated.

All models are compared on identical forecast timestamps.
State-space forecasts are generated sequentially, with each
realized outcome entering the filter only after the
corresponding forecast has been produced.

Fixed HAR-X and TVP-HAR-X use the same predictor set, so their
performance difference reflects coefficient variation rather
than unequal observable information. Forecast-loss comparisons
also use HAC standard errors to account for serial dependence
and heteroskedasticity.

These procedures reduce the scope for data leakage and ensure
that the reported loss differences arise from the model
specifications evaluated under common conditions.

\subsection{Limitations}

\subsubsection{Single Market and Contract}

The analysis is based on the Binance BTCUSDT USD-margined
perpetual futures contract. The results may differ for Bitcoin
spot markets, other exchanges, alternative cryptocurrency
assets, or less liquid contracts. Exchange-specific trading
mechanisms and participant composition may affect both market
activity and volatility dynamics.

\subsubsection{Single Forecast Horizon}

The study considers only one-hour-ahead volatility. Predictor
importance and model rankings may change at shorter or longer
horizons. Order-flow variables may be more informative over
several minutes, while longer-horizon volatility components
may contribute more to multi-hour or daily forecasts.

\subsubsection{Limited Microstructure Information}

The available data include candlestick volume, trade count, and
taker-buy volume, but exclude the complete limit order book.
Bid--ask spreads, market depth, cancellations, liquidations,
price impact, open interest, and funding conditions are not
observed directly.

The market microstructure variables used here should therefore
be interpreted as aggregate high-frequency trading proxies.

\subsubsection{Linear Predictor Effects}

Both HAR-X specifications impose linear effects in the
observation equation. The relationship between activity and
future volatility may be nonlinear, particularly during
extreme market conditions.

A signed aggressive-buying measure may also understate the
effect of order-flow imbalance when both strong buying and
strong selling are associated with elevated volatility.
Nonlinear and absolute-imbalance terms would require
pre-specification and evaluation on a new holdout sample.

\subsubsection{Restricted State Covariance}

All seven time-varying slope coefficients share the common
state innovation variance \(q_{\beta}\). This restriction
improves parsimony and numerical stability but imposes a common
prior rate of variation across predictor effects.

Separate state variances could allow trading activity,
short-horizon volatility, and weekly volatility coefficients
to evolve at different rates. Such a specification would add
hyperparameters and may increase estimation uncertainty.

\subsubsection{Retransformation from Log Variance}

The models forecast log realized variance, and the variance
forecast used in QLIKE is obtained as

\[
\widehat{RV}_t
=
\exp
\left(
\widehat{y}_t
\right).
\]

In general,

\[
E
\left[
RV_t
\mid
\mathcal{F}_{t-1}
\right]
\neq
\exp
\left(
E
\left[
\log RV_t
\mid
\mathcal{F}_{t-1}
\right]
\right).
\]

The direct exponential transformation is closer to a
conditional median than to a conditional mean when log
forecast errors have positive variance. The same
transformation is applied to every model, preserving internal
comparability. Future work could evaluate lognormal or
nonparametric bias corrections for variance-scale forecasting.

\subsubsection{Seasonality and Long-Memory Structure}

The feature set contains six-hour, daily, and weekly volatility
averages but no explicit hour-of-day or day-of-week indicators.
Continuous cryptocurrency trading may still display seasonal
patterns associated with participation across geographic
regions.

The significant long-horizon Ljung--Box statistics also
suggest that additional seasonal or long-memory dependence
remains.

\subsection{Directions for Further Research}

The analysis can be extended in several directions. Applying
the same design to multiple contracts and exchanges would show
whether the predictive value of trading activity generalizes
across markets. Evaluating several forecast horizons would
clarify how the importance of HAR components and activity
variables changes with the prediction interval.

Richer high-frequency predictors could include order-book
depth, bid--ask spreads, liquidation volume, open interest, and
funding rates. Nonlinear transformations of trade count and
order-flow imbalance may also capture effects missed by the
linear specification.

Alternative state dynamics provide another extension.
Mean-reverting coefficients or regime-switching models may be
appropriate when forecasting relationships alternate among
recurring market states. Forecast combinations could also
exploit the relative stability of Fixed HAR-X and the adaptive
behavior of TVP-HAR-X.

Any extension informed by the present test results would
require a new validation period and an untouched final holdout
sample.